
\section{maincontent}

\begin{frame}
\frametitle{Magnetic-Field Effect in High-Order Above-Threshold Ionization \\
\small (Phys. Rev. Lett. 128, 023201 (2022))}
\framesubtitle{Surrounding concepts}

\small
\begin{textblock*}{100pt}(30pt,60pt)
\begin{itemize}
\item Polarization

\includegraphics[scale=.75]{images/polcls.png}
\item Nondipole effects
\begin{align}
\footnotesize
H = \frac{1}{2}
\Big( \boldsymbol{p} + q\boldsymbol{A}(x,t) \Big)^2 + V_0, \nonumber \\
\boldsymbol{A}(x,t) \sim A(t) + \frac{x}{c}F(t) 
- \frac{1}{2}\frac{x^2}{c^2}\frac{d}{dt}F(t). \nonumber
\end{align}
\item Relevant equations
\begin{align}
\footnotesize
\boldsymbol{F} = -\frac{\partial \boldsymbol{A}}{\partial t}, \hspace{10pt} 
\nabla \times \boldsymbol{B} = 
   \frac{1}{c^2}\frac{\partial\boldsymbol{F}}{\partial t}.
 \nonumber
\end{align}
\end{itemize}
\end{textblock*}


\small
\begin{textblock*}{200pt}(230pt,60pt)
\begin{itemize}
\item Ponderomotive energy
\begin{align}
\footnotesize
U_p = \frac{F_0^2}{4\omega^2}. \nonumber
\end{align}
\item Cross-section $\sigma$ \\
{\scriptsize Quantity proportional to the rate at which a particular 
radiation-target interaction occurs.  
If the incoming radiation is considered as being composed 
by quanta, then the cross section is a scattering rate.}
\end{itemize}
\end{textblock*}

%\begin{textblock*}{400pt}(30pt,130pt)
%\centering\footnotesize
%\begin{align}
%H = \frac{1}{2}
%\Big( \boldsymbol{p} + \boldsymbol{A}(x,t) \Big) \nonumber \\
%\boldsymbol{A}(x,t) \sim A(t) + \frac{x}{c}E(t) 
%- \frac{1}{2}\frac{x^2}{c^2}\frac{d}{dt}E(t) \nonumber
%\end{align}
%
%\end{textblock*}

\end{frame}



\begin{frame}
\frametitle{Magnetic-Field Effect in High-Order Above-Threshold Ionization \\
\small (Phys. Rev. Lett. 128, 023201 (2022))}
\framesubtitle{The state of the problem}

\small
\begin{textblock*}{400pt}(30pt,60pt)
\begin{itemize}
\item Conservation of momentum
\begin{align}
\footnotesize
\left\langle p_{x,e} \right\rangle + \left\langle p_{x,\rm ion} \right\rangle
 = \frac{E_\gamma}{c} 
 = \frac{E_e}{c} + \frac{I_p}{c}. \nonumber
\end{align}
\item How is the photon energy shared between electrons and ions?
\begin{itemize}
 \footnotesize
 \item[$\bullet$] For high photon energies, \hspace{10pt}
   $\left\langle p_{x,e}\right\rangle = 8/5\times E_e/c.$
 \item[$\bullet$] For recollision free strong field ionization, %\footnote{\Tiny Chelkowski \textit{et al.}, Phys. Rev. Lett. 113, 263005, (2014)}
 \hspace{10pt}
   $\left\langle p_{x,e} \right\rangle \sim E_e/c + 1/3\times I_p.$
 \item[$\bullet$] For linearly polarized fields, around $2U_p$ region, 
   \hspace{10pt}
   $\left\langle p_{x,e}\right\rangle \sim E_e/c.$
\end{itemize}
\item For rescattered electrons whose energy can reach up to $10U_p$, 
the question is open. \\
Still, Brennecke \textit{et al.} predicted that for $\rm He$,
$\langle p_{x,e}\rangle$ levels off and remains constant.
\item Some interesting comments
\begin{itemize}
 \footnotesize
% \item[$\bullet$] Recollision processes are at the heart of phenomema such as 
%HHG and laser-induced electron diffraction
 \item[$\bullet$] Requirements of high intensities and \underline{long} 
wavelengths place HHG and LIED outside of the dipole oasis.
 \item[$\bullet$] Nondipole effects have been deepened by the analysis 
of NSI, field inhomogeneity, nonadiabaticity and weak coulomb effects...
\end{itemize}

\end{itemize}
\end{textblock*}

%\begin{textblock*}{400pt}(30pt,130pt)
%\centering\footnotesize
%\begin{align}
%H = \frac{1}{2}
%\Big( \boldsymbol{p} + \boldsymbol{A}(x,t) \Big) \nonumber \\
%\boldsymbol{A}(x,t) \sim A(t) + \frac{x}{c}E(t) 
%- \frac{1}{2}\frac{x^2}{c^2}\frac{d}{dt}E(t) \nonumber
%\end{align}
%
%\end{textblock*}

\end{frame}


\begin{frame}
\frametitle{Magnetic-Field Effect in High-Order Above-Threshold Ionization \\
\small (Phys. Rev. Lett. 128, 023201 (2022))}
\framesubtitle{What to expect from this paper}

\small
\begin{textblock*}{400pt}(30pt,60pt)
\begin{itemize}
\item First experimental results on the photon-momentum transfer
for the ionization of $\rm Xe$ atoms by linearly polarized light 
beyond the classical cutoff.
\item First experimental test of the Brennecke proposition, 
(which was about $\rm He$).
\vspace{10pt}
\item Quick experimental observations
\begin{itemize}
\item[$\bullet$] A breakdown of $\left\langle p_{x,e} \right\rangle$
is found as soon as energy gets higher than $2U_p$.
\item[$\bullet$] Still, the pleateau-like behaviour predicted by Brennecke 
for $\rm He$ is not seen.
\item[$\bullet$] Chances are that the momentum transfer is target sensitive.
\item[$\bullet$] Present paper results are for $\rm Xe$. 
\end{itemize}
\end{itemize}

\end{textblock*}

\end{frame}


\begin{frame}
\frametitle{Magnetic-Field Effect in High-Order Above-Threshold Ionization \\
\small (Phys. Rev. Lett. 128, 023201 (2022))}
\framesubtitle{Experimental setup}

\small

\begin{textblock*}{100pt}(280pt,70pt)
\includegraphics[scale=.30]{images/medium1.png}
\end{textblock*}

\begin{textblock*}{250pt}(30pt,60pt)
\begin{itemize}
\footnotesize
\item Linear-photon-momentum transfer $\left\langle \Delta p_x\right\rangle$
is around $4\times 10^{-4}\,\rm a.u$ for Ti:Sapphire  $(800\,\rm nm)$.
\item Making use of the following strategy
\begin{itemize}
\footnotesize
\item[$\bullet$] $(800\,\rm nm$,\ 10\,\rm kHz$,\ 25\,\rm fs)$, 
Coherent Legend Elite Duo \\
\Tiny
SANTA CLARA, Calif., July 1, 2010 --  A new Ti:sapphire ultrafast laser amplifier from Coherent Inc., the Legend Elite Duo-HP, offers >15 W of output power and pulse energy of up to 10 mJ. The kilohertz-class thermoelectrically cooled amplifier can be configured to provide output pulse durations of <25 fs, <35 fs, <130 fs or 0.5 to 2 ps with repetition rates ranging from 1 to 10 kHz.
\footnotesize
\item[$\bullet$] Split, then focused back from 2 opposite directions 
on the same spot of supersonic $\rm Xe$ gas.
\item[$\bullet$] In the reaction volume, 
$7\times 10^{13}\,{\rm W/cm^2}\;\; \left(U_p = 4.2\,\rm eV\right)$.
\item[$\bullet$] Electrons and ions guided to opposite directions 
with a static field of
$29.8\,\rm V/cm$.
\item[$\bullet$] Symmetry of the experiment along the polarization 
direction ($z$-axis) prompted the symmetrisation 
of the data in regard that same direction.
\end{itemize}
\end{itemize}
\end{textblock*}
\end{frame}


\begin{frame}
\frametitle{Magnetic-Field Effect in High-Order Above-Threshold Ionization \\
\small (Phys. Rev. Lett. 128, 023201 (2022))}
\framesubtitle{Two-dimensional PEMD along the polarization $z-$ and 
light propagation $x-$ axis}

\begin{textblock*}{100pt}(300pt,70pt)
\includegraphics[scale=.30]{images/medium2.png}
\end{textblock*}

\small
\begin{textblock*}{250pt}(30pt,60pt)
\begin{itemize}
\footnotesize
\item Figure 2(a) displays the $2D$ distribution of the particules momentum.
\begin{itemize}
\footnotesize
\item[$\bullet$] $|p_z| \leqslant 0.2\,\rm a.u.$ electrons experience strong
Coulomb focusing.
\item[$\bullet$] $ 0.2\,{\rm a.u.} \leqslant |p_z| \leqslant 2\sqrt{U_p}$,
distribution is dominated by photon holography and ATI rings are seen.
\item[$\bullet$] $ 2\sqrt{U_p} \leqslant |p_z| $, distribution dominated by 
HATI. \\
Over the $10U_p$, BRR can be identified.
\end{itemize}
\item Figure 2(b) shows the distribution (or the number) of electron
having a precise momentum regarding the polarization direction. 
\item Figure 2(c) shows the distribution of transferred momentum $p_x$ along 
the light propagation direction for each $p_z$.
\begin{itemize}
\footnotesize
\item[$\bullet$] For the intermediate energies, overall agreement.
\item[$\bullet$] For $p_z < 2\sqrt{U_p}$, experimental results shift
negative values.
\item[$\bullet$] For $p_z > 2\sqrt{U_p}$, values deviate from the expected 
shape. 
\end{itemize}
\end{itemize}
\end{textblock*}



\end{frame}



\begin{frame}
\frametitle{Magnetic-Field Effect in High-Order Above-Threshold Ionization \\
\small (Phys. Rev. Lett. 128, 023201 (2022))}
\framesubtitle{Modelling and theoretical calculations}

%\begin{textblock*}{100pt}(280pt,70pt)
%\includegraphics[scale=.30]{images/medium3.png}
%\end{textblock*}

\small
\begin{textblock*}{400pt}(30pt,60pt)
\begin{itemize}
\footnotesize
\item Explanation of the local minimum structure.
\item Simple backreflection model (BRM) based on some 
nondipole generalization of the 3 steps model.
\begin{itemize}
\footnotesize
\item[$\bullet$] First steps are the exact same.
\item[$\bullet$] Incident velocity is flipped around by the magnetic field 
and augmented by $v_{{\rm in},x} = v_{{\rm in},z}^2/(2c)$ 
along the light propagation direction.
\item[$\bullet$] For $\rm Xe^+$, the differential electron scattering 
cross-section has a local maximum at $180^{\degree}$ (backscattering) 
which induces the maximum in the $p_x$-direction near the $p_z$-axis.
\item[$\bullet$] They therefore \underline{estimate} the nondipole shift
by "assuming" that all electrons are backscattered.
\item[$\bullet$] Recollision slows down the electron momentum along $p_x$-, 
while it is boosted along $p_z$-direction.
\item[$\bullet$] Two trajectories born per optical cycle 
contribute for the same $p_z$, termed "long-" and "short-".
\end{itemize}
\item Low Frequency Approximation (LFA) is implemented in order to complete 
the above model.
\begin{itemize}
\footnotesize
\item[$\bullet$] Based on complex-valued electron trajectories.
\item[$\bullet$] Accurately incorporate the field free differential. 
\item[$\bullet$] 
\begin{align}
M(\boldsymbol{p}) = P(\boldsymbol{p}) 
T\left[ \boldsymbol{v}_{\rm in}(\boldsymbol{p}), 
\boldsymbol{v}_{\rm out}(\boldsymbol{p})\right], 
\hspace{20pt}\sigma \sim \left|  T\right|^2. \nonumber
\end{align}
\end{itemize}

\end{itemize}
\end{textblock*}

\end{frame}


\begin{frame}
\frametitle{Magnetic-Field Effect in High-Order Above-Threshold Ionization \\
\small (Phys. Rev. Lett. 128, 023201 (2022))}
\framesubtitle{Result of the classical and quantum-orbit analysis}

\begin{textblock*}{100pt}(330pt,70pt)
\includegraphics[scale=.30]{images/medium3.png}
\end{textblock*}

\small
\begin{textblock*}{300pt}(30pt,60pt)
\begin{itemize}
\footnotesize
\item (a) Momentum-shift estimate from classical BRM and LFA.
\begin{itemize}
\footnotesize
\item[$\bullet$] Nondipole shit in LFA long- and short- trajectories 
similar to the BRM.
\item[$\bullet$] Nonetheless, due to nonadiabatic initial velocities, 
difference between short- and long- orbits is larger.
\end{itemize}
\item (b) Probability of quantum orbits with final momentum $p_x = 0$ in LFA.
\begin{itemize}
\footnotesize
\item[$\bullet$] Though both orbits contribute, the shift is strongly 
influenced by the relative weight between them.
\item[$\bullet$] Relative weight mainly determined by the energy-dependent
cross-section.
\item[$\bullet$] Long- dominates for $p_z \leqslant 1.3\,\rm a.u$, as 
short- becomes comparable for $1.3\,{\rm a.u.} 
\leqslant p_z \leqslant 1.7\,{\rm a.u.}$
\item[$\bullet$] Transition from long- to short-, then back to long trajectory 
causing the minimum structure in the shift $p_x$.
\end{itemize}
\item (c) Focal volume-averaged LFA result calculated for 
the Tong-Lin potential.
\begin{itemize}
\footnotesize
\item[$\bullet$] Reproduces the experimentally observed minimum structure.
\item[$\bullet$] Confirms the prediction for $p_z \geqslant 1.6\,\rm a.u$.
\end{itemize}
\end{itemize}
\end{textblock*}


\end{frame}




\begin{frame}
\frametitle{Magnetic-Field Effect in High-Order Above-Threshold Ionization \\
\small (Phys. Rev. Lett. 128, 023201 (2022))}
\framesubtitle{Model potential dependence of the nondipole shift}

\begin{textblock*}{100pt}(280pt,70pt)
\includegraphics[scale=.30]{images/medium4.png}
\end{textblock*}

\small
\begin{textblock*}{250pt}(30pt,60pt)
\begin{itemize}
\footnotesize
\item Explanation of the difference between $\rm He$ and $\rm Xe$.
\vspace{10pt}
\item Model potential $V(r)$ expression
\begin{enumerate}
\footnotesize
\item
$-\left( Z + a_1 e^{-a_2 r} + a_3 e^{-a_4 r} + a_5 e^{-a_6 r} \right)/r$
(Tong-Lin~),
\item
$-2\left( N\Gamma - Z \right)/r, \;\; \mbox{with}\;\;
\Gamma = 1 - \left[ \left( e^{r/d}  \right)B + 1 \right]$ (GSZ),
\item
$-\left( Z + a_1 e^{-a_2 r} \right)/r$ (effective).
\end{enumerate}
\item Predicted shifts different in the  HATI region, which implies 
that the electron-ion interaction plays an important role.
\item Tong-Lin potential (1) gives the most accurate results.
\item Comparison shows that shifts are very sensitive to the details 
of the scattering process.
\item This also explains why results are different from $\rm He$, since 
the specificities of the target are very important for the dynamics.
\item Good agreement between theory and prediction 
for $p_z \leqslant 1.0\,\rm a.u.$
\end{itemize}
\end{textblock*}



\end{frame}


\begin{frame}
\frametitle{Magnetic-Field Effect in High-Order Above-Threshold Ionization \\
\small (Phys. Rev. Lett. 128, 023201 (2022))}
\framesubtitle{Closure}

\small
\begin{textblock*}{400pt}(30pt,60pt)
\begin{itemize}
\footnotesize
\item Partitioning of the linear photon momentum between 
electron and nucleus altered by the rescattering process. 
\item Depending on the electron energy compared to $2U_p$, the share 
taken by the electron drastically different.
\item In the case of $\rm Xe$, domination of the distribution by 
backscattered electrons in the  $2\sqrt{U_p} \leqslant p_z$ region.
\item Recollision with the parent ion transfers 
most of the momentum to the core, which halves the momentum altogether,
especially in the $p_x$ direction. 

\vspace{20pt}

\item Sensitivity of the nondipole shift to the inner structure 
of the parent ion turns it into a promising tool 
for atomic and molecular imaging.
\end{itemize}
\end{textblock*}


%\begin{textblock*}{100pt}(260pt,70pt)
%\includegraphics[scale=.30]{images/medium1.png}
%\end{textblock*}

\end{frame}
